\chapter{Budget Constraints and Consumer Choice}
\label{chap:consumer-choice}

\section{Learning objectives}

After completing this chapter, you should be able to:
\begin{itemize}
  \item write and graph a consumer's budget set and budget line;
  \item explain the effects of income and price changes on the budget constraint;
  \item identify utility-maximising bundles graphically and algebraically; and
  \item distinguish interior from corner solutions.
\end{itemize}

\section{Preferences meet constraints}

Preferences describe what a consumer would like; a budget constraint describes what the consumer can afford. Consumer choice combines the two. The consumer selects the most preferred bundle among all bundles in the feasible set.

Let income be $I$, prices be $p_x$ and $p_y$, and quantities be $x$ and $y$. The \keyterm{budget constraint} is
\[
p_xx+p_yy\leq I.
\]
The left-hand side is total expenditure. The inequality allows the consumer to spend less than income. With monotonic preferences, the consumer normally chooses a bundle on the \keyterm{budget line}, where all income is spent:
\[
p_xx+p_yy=I.
\]

\section{The budget line}

Solving the budget line for $y$ gives
\[
y=\frac{I}{p_y}-\frac{p_x}{p_y}x.
\]
The vertical intercept is $I/p_y$: the maximum quantity of $y$ if all income is spent on $y$. The horizontal intercept is $I/p_x$. The slope is $-p_x/p_y$, the negative of the relative price of $x$ in terms of $y$.

For example, suppose $I=100$, $p_x=10$, and $p_y=5$. Then
\[
10x+5y=100,
\qquad\text{or}\qquad
y=20-2x.
\]
The intercepts are $(0,20)$ and $(10,0)$, and the slope is $-2$. One extra unit of $x$ costs two units of $y$: this is the opportunity cost of $x$ in terms of $y$.

\section{Changes in income and prices}

When income changes but prices remain fixed, the budget line shifts in a parallel way. A higher income moves both intercepts outward; a lower income moves both inward.

When the price of one good changes, the budget line pivots around the intercept of the good whose price did not change. If $p_x$ falls, $I/p_x$ rises, so the horizontal intercept moves outward while the vertical intercept stays fixed. The line becomes flatter because $x$ is relatively cheaper.

These changes are fundamental for demand analysis. They show both a change in purchasing power and a change in the relative attractiveness of goods.

\section{Utility maximisation graphically}

With smooth, monotonic, convex preferences, the optimal bundle is normally the point where the highest attainable indifference curve is tangent to the budget line. At this interior optimum,
\[
MRS_{x,y}=\frac{p_x}{p_y}.
\]
The consumer's willingness to trade $y$ for $x$ equals the trade-off required by the market. If the MRS were greater than the price ratio, the consumer would be willing to give up more $y$ for $x$ than the market requires, so moving toward more $x$ would raise utility. If the MRS were lower, moving toward more $y$ would raise utility.

\section{Utility maximisation algebraically}

Consider $u(x,y)=xy$, $I=100$, $p_x=10$, and $p_y=5$. From the budget constraint,
\[
y=20-2x.
\]
Substitute into utility:
\[
u(x)=x(20-2x)=20x-2x^2.
\]
The first derivative is $20-4x$. Setting it equal to zero gives $x^*=5$. Substitution into the budget line gives $y^*=10$. The bundle $(5,10)$ exhausts income and maximises utility.

The same answer follows from the tangency condition. Here $MRS_{x,y}=y/x$, while $p_x/p_y=2$. Setting $y/x=2$ gives $y=2x$. Combined with $10x+5y=100$, this yields $x=5$ and $y=10$.

\section{Interior and corner solutions}

The tangency condition applies only to an \keyterm{interior solution}, where positive quantities of both goods are consumed. With perfect substitutes, the best bundle may be a \keyterm{corner solution}: all income is spent on one good.

Suppose $u(x,y)=x+2y$, $p_x=4$, and $p_y=1$. One euro spent on $x$ produces $1/4$ utility units, while one euro spent on $y$ produces $2$ utility units. The consumer spends all income on $y$. There is no interior tangency because the MRS is constant and does not equal the price ratio.

For perfect complements, the consumer chooses the kink on the highest attainable indifference curve. The budget line need not be tangent to the indifference curve.

\section{From consumer choice to demand}

Solving the consumer's problem for different prices produces an individual demand function. The next chapter develops this link in detail. For now, remember that demand is not simply a psychological desire for a good; it is the quantity that maximises utility given prices, income, and preferences.

\section*{Chapter Review}

\subsection*{Key concepts}
\begin{itemize}
  \item \textbf{Budget set}: all affordable bundles, $p_xx+p_yy\leq I$.
  \item \textbf{Budget line}: bundles that exactly exhaust income.
  \item \textbf{Relative price}: $p_x/p_y$, the opportunity cost of $x$ in units of $y$.
  \item \textbf{Interior solution}: an optimum involving positive amounts of both goods.
  \item \textbf{Corner solution}: an optimum where at least one good is consumed in zero quantity.
\end{itemize}

\subsection*{Key equations}
\[
p_xx+p_yy=I,
\qquad
y=\frac{I}{p_y}-\frac{p_x}{p_y}x,
\qquad
MRS_{x,y}=\frac{p_x}{p_y}.
\]

\section*{Exercises}

\subsection*{Conceptual questions}
\begin{enumerate}[label=\arabic*.]
  \item Distinguish the budget set from the budget line.
  \item Explain why a consumer with monotonic preferences normally exhausts their budget.
  \item What is the economic interpretation of the slope of a budget line?
  \item Explain the difference between an income change and a change in the price of one good.
  \item Why does the tangency condition not apply to all consumer-choice problems?
\end{enumerate}

\subsection*{Graphical exercises}
\begin{enumerate}[resume,label=\arabic*.]
  \item Let $I=120$, $p_x=12$, and $p_y=6$. Draw the budget line and label its intercepts and slope.
  \item Starting from Exercise 6, show graphically what happens when income rises to $180$.
  \item Starting from Exercise 6, show graphically what happens when $p_x$ falls from $12$ to $6$.
\end{enumerate}

\subsection*{Analytical exercises}
\begin{enumerate}[resume,label=\arabic*.]
  \item A consumer has $I=200$, $p_x=20$, $p_y=10$, and $u(x,y)=xy$. Solve for the optimal bundle by substitution and verify it using $MRS_{x,y}=p_x/p_y$.
  \item Solve the consumer problem when $I=100$, $p_x=10$, $p_y=5$, and $u(x,y)=x^{2/3}y^{1/3}$.
  \item Let $u(x,y)=x+3y$, $p_x=2$, $p_y=1$, and $I=40$. Find the optimal bundle and explain why it is a corner solution.
\end{enumerate}

\subsection*{Exam-style multiple-choice questions}
\begin{enumerate}[resume,label=\arabic*.]
  \item The slope of the budget line $p_xx+p_yy=I$ is: (a) $p_x/p_y$; (b) $-p_x/p_y$; (c) $I/p_x$; or (d) $I/p_y$.
  \item If income rises while prices are unchanged, the budget line: (a) pivots around the horizontal intercept; (b) becomes steeper; (c) shifts outward in parallel; or (d) does not change.
  \item At an interior optimum, the consumer chooses a bundle where: (a) income is zero; (b) $MRS_{x,y}=p_x/p_y$; (c) marginal utility is zero for both goods; or (d) both goods have the same price.
\end{enumerate}

\subsection*{PPE discussion question}

\noindent A budget constraint is often treated as an individual fact. Discuss how public policy and institutions can shape it, considering taxes, public services, discrimination, access to credit, and social insurance.
